% !TEX root = ../main.tex

A jelen fejezetben bemutatom a CDR adatfeldolgozó rendszer tervezési döntéseit: milyen technológiai alternatívákat mérlegeltem, és az adott korlátok mellett miért az alábbi eszközöket választottam. A rendszer célja a YJMob100K adathalom (5.\ fejezet) feldolgozása, a 4.\ fejezetben ismertetett mobilitási metrikák kiszámítása, valamint a származtatott mutatók GIS alapú megjelenítése.

\subsection{Tervezési szempontok és korlátok}

A rendszerterv kialakítását három fő szempont határozta meg:
\begin{enumerate}
    \item \textbf{Adatméret:} A YJMob100K adathalom két fő CSV fájlja összesen mintegy 141~millió sort tartalmaz. Az adatbázis-motornak képesnek kell lennie ennyi rekord hatékony kezelésére anélkül, hogy a teljes adatot a memóriába kellene tölteni.
    \item \textbf{Hardverkorlátok:} A fejlesztő gép 16~GB RAM-mal és 6 magos Intel i5-8400 processzorral rendelkezik. Elosztott cluster nem áll rendelkezésre, így a választott eszközöknek egyetlen gépen is hatékonynak kell lenniük.
    \item \textbf{Reprodukálhatóság:} A teljes pipeline-nak szkriptelhetőnek, verziókezelhetőnek és parancssori futtatással újrafuttathatónak kell lennie, külső szerver vagy GUI nélkül.
\end{enumerate}

\subsection{Technológiai alternatívák értékelése}

Az alábbiakban áttekintem az egyes funkcióterületek mérlegelt alternatíváit, és megindokolom a végső választást.

\subsubsection{Adatbázis-motor}

A CDR feldolgozás irodalmában a relációs adatbázisok (PostgreSQL, Microsoft SQL Server), az elosztott big-data rendszerek (Apache Spark + Hadoop HDFS), valamint az újabb beágyazott OLAP megoldások (DuckDB, SQLite) fordulnak elő \cite{zhang2019connected, pinter2020artificial}.

\begin{itemize}
    \item \textbf{PostgreSQL / PostGIS} -- széles körben használt nyílt forráskódú relációs adatbázis, amely a PostGIS kiterjesztéssel natív térbeli indexelést és GIS-funkciókat nyújt. Hátránya, hogy külön szerverfolyamatot igényel, és a 141 millió soros OLAP jellegű aggregációknál (teljes tábla átlagolása, GROUP BY a teljes adaton) sor-orientált tárolása miatt lassabb, mint egy oszlopos motor.
    \item \textbf{Apache Spark} -- elosztott számítási keretrendszer, amely a szakirodalomban gyakori választás CDR feldolgozásra \cite{zhang2019connected}. Legnagyobb előnye a horizontális skálázhatóság, azonban egyetlen gépen a JVM-alapú futtatókörnyezet magas memóriaigénye és a konfigurációs komplexitás nem indokolt.
    \item \textbf{DuckDB} -- beágyazott (in-process), oszlopos OLAP adatbázis-motor, amely közvetlenül a Python folyamatba integrálódik. Nem igényel szervertelepítést, a CSV fájlokat streaming módban olvassa (nem kell a teljes adat a memóriában), és a vektorizált végrehajtó motorja révén a százmilliós aggregációk is másodpercek alatt futnak 16~GB RAM mellett.
\end{itemize}

A \textbf{DuckDB} mellett döntöttem, mivel az adathalom ugyan nagy, de egyetlen gépre fér; az oszlopos tárolás és a vektorizált végrehajtás ideális a mobilitási metrikák GROUP BY-intenzív SQL lekérdezéseihez; és a Python API-n keresztül közvetlenül pandas DataFrame-ekké alakíthatók az eredmények. Az SQL nyelv használata azt is biztosítja, hogy a lekérdezések a jövőben más motoron (pl.\ PostgreSQL) is futtathatók maradjanak.

\subsubsection{Programozási nyelv és ETL megközelítés}

A feldolgozó pipeline nyelveként \textbf{Python} volt a természetes választás, mivel az adatelemzési ökoszisztéma (pandas, scikit-learn, matplotlib) egységes keretet biztosít az ETL, a gépi tanulás és a vizualizáció számára. Alternatívaként felmerülő R-nyelv GIS-támogatása kevésbé kiforrott, az SSIS (SQL Server Integration Services) pedig zárt forráskódú és platformfüggő.

Az ETL logikát közvetlenül Python szkriptekbe szerveztem, modul-szintű felelősségmegosztással. Egy-egy szkript egy-egy pipeline-fázist valósít meg, és a parancssorból futtatható. A köztes eredményeket Parquet formátumban is exportálom, ami hatékony oszlopos tárolást biztosít és biztonsági mentésként is szolgál.

\subsubsection{Térbeli feldolgozás (GIS)}

A GIS komponens feladata a YJMob100K rácskoordinátáinak georeferálása és térképes megjelenítése. A mérlegelt alternatívák:
\begin{itemize}
    \item \textbf{QGIS} -- asztali GIS alkalmazás, amely manuális rétegkezelést és interaktív térképszerkesztést kínál. Hátránya, hogy a munkafolyamat nehezen szkriptelhető és nem reprodukálható automatikusan.
    \item \textbf{PostGIS} -- a PostgreSQL térbeli kiterjesztése; erős a térbeli lekérdezésekben, de külön adatbázis-szervert igényel, és a vizualizációhoz további eszköz (pl.\ QGIS) szükséges.
    \item \textbf{GeoPandas + pyproj + contextily} -- Python könyvtárcsalád, amely programozottan végzi a koordináta-transzformációkat, GeoDataFrame-eket kezel, és OpenStreetMap alaptérképekre rajzol. Teljes mértékben szkriptelhető és a Jupyter Notebook-okba integrálható.
\end{itemize}

A \textbf{GeoPandas + pyproj} kombinációt választottam: a pyproj a koordináta-transzformációt végzi (EPSG:2449 $\rightarrow$ WGS84), a GeoPandas a térbeli adatszerkezeteket kezeli, a contextily pedig az OpenStreetMap csemperéteget biztosítja alaptérképként. Ez a megoldás lehetővé teszi, hogy a teljes GIS munkafolyamat Python kódból futtatható és reprodukálható legyen.

\subsubsection{Anomáliadetekció}

A rendszerterv része a viselkedési anomáliák felismerése is, amelyet az adatkészlet második részének vészhelyzeti időszaka (60--74.\ nap) tesz lehetővé. A scikit-learn könyvtár Isolation Forest algoritmusát választottam, mivel felügyelet nélküli módszer (nem igényel címkézett adatot), jól skálázódik a többszázezres mintaméretre, és a Python ökoszisztéma részét képezi.

\subsubsection{Vizualizáció}

A vizualizációhoz a \textbf{matplotlib + seaborn} párost választottam a statikus, tudományos minőségű ábrák elkészítéséhez. Az interaktív megoldásokat (Plotly, Folium, kepler.gl) mérlegeltem, de a dolgozat célja nyomtatott minőségű ábrák előállítása, amelyhez a matplotlib jobban illeszkedik. A kontextushoz az OpenStreetMap alaptérképeket a contextily könyvtár biztosítja.

\subsection{A választott technológiai stack összefoglalása}

Az alábbi táblázat foglalja össze az egyes rétegek kiválasztott eszközeit:

\begin{table}[H]
\centering
\caption{A megvalósított rendszer technológiai stack-je.}
\label{tab:tech_stack}
\begin{tabular}{l l l}
\hline
\textbf{Réteg} & \textbf{Eszköz} & \textbf{Szerep} \\
\hline
Adatbázis       & DuckDB            & Oszlopos OLAP, in-process \\
ETL nyelv       & Python            & Pipeline szkriptek \\
Adathíd         & pandas            & Python $\leftrightarrow$ DuckDB \\
Georeferálás    & pyproj            & EPSG:2449 $\rightarrow$ WGS84 \\
Térbeli adat    & GeoPandas         & GeoDataFrame kezelés \\
Alaptérkép      & contextily        & OpenStreetMap csempék \\
Anomáliadetekció & scikit-learn     & Isolation Forest \\
Ábrák           & matplotlib, seaborn & Statikus vizualizáció \\
Köztes formátum & Parquet           & Oszlopos export/mentés \\
\hline
\end{tabular}
\end{table}


\subsection{Magasszintű architektúra}

A rendszer négy fő rétegből áll (\autoref{fig:arch}):

\begin{enumerate}
    \item \textbf{Adatforrás réteg.} A nyers CSV fájlok (YJMob100K dataset1, dataset2) és a kiegészítő adatok (POI kategóriák). Ezeket a DuckDB \texttt{read\_csv} streaming olvasója dolgozza fel, tehát a teljes CSV soha nem kerül egyben a memóriába.

    \item \textbf{Perzisztens tároló (DuckDB).} A betöltött adatok DuckDB táblákba kerülnek, amelyek oszlopos formátumban, tömörítve tárolódnak egyetlen fájlban (\texttt{mobility.duckdb}). Ez egyben a Data Warehouse szerepet is betölti, mivel a dimenziós táblákat is itt hozom létre.

    \item \textbf{Analitikai réteg.} A mobilitási metrikákat SQL CTE-láncokon (Common Table Expression) keresztül számítom ki a DuckDB-ben, majd az eredményeket Parquet fájlokba is exportálom. Az anomáliadetekciót a scikit-learn végzi Python-ban, a DuckDB-ből kiolvasott feature mátrixon.

    \item \textbf{Megjelenítési réteg.} A származtatott mutatókat a GeoPandas georeferálja, a matplotlib + contextily pedig térképes és statisztikai ábrákon jeleníti meg.
\end{enumerate}

\begin{figure}[H]
\centering
\begin{tikzpicture}[
    scale=0.70,
    transform shape,
    node distance=0.8cm,
    auto,
    block/.style={
        rectangle,
        draw=black!60,
        fill=blue!5,
        thick,
        text width=3.0cm,
        align=center,
        font=\small,
        rounded corners=2pt,
        minimum height=3.2em,
        drop shadow
    },
    arrow/.style={
        -{Latex[length=2mm]},
        thick,
        draw=black!70
    }
]
    \node [block, fill=orange!10] (step1) {\textbf{CSV fájlok}\\ \footnotesize(YJMob100K)};
    \node [block, right=of step1, fill=green!10] (step2) {\textbf{DuckDB}\\ \footnotesize(OLAP tároló \\ + csillagséma)};
    \node [block, right=of step2, fill=cyan!10] (step3) {\textbf{Metrikák}\\ \footnotesize(SQL CTE-lánc \\ + Parquet)};
    \node [block, right=of step3, fill=yellow!20] (step4) {\textbf{Anomália}\\ \footnotesize(Isolation \\ Forest)};
    \node [block, below=0.7cm of step3, fill=purple!10] (step5) {\textbf{GIS + ábrák}\\ \footnotesize(geo.py, matplotlib)};

    \draw [arrow] (step1) -- node[above, font=\scriptsize]{ingest.py} (step2);
    \draw [arrow] (step2) -- node[above, font=\scriptsize]{features.py} node[below, font=\scriptsize]{schema.py} (step3);
    \draw [arrow] (step3) -- node[above, font=\scriptsize]{anomaly.py} (step4);
    \draw [arrow] (step2) |- (step5);
    \draw [arrow] (step3) -- (step5);
    \draw [arrow] (step4) -| (step5);
\end{tikzpicture}
\caption{A megvalósított CDR feldolgozó rendszer magas szintű architektúrája. A négy ETL-fázis (ingest $\to$ features/schema $\to$ anomaly) a DuckDB tárolón keresztül kommunikál; a megjelenítési réteg (\texttt{geo.py} + matplotlib) közvetlenül olvassa a DuckDB-t és a Parquet outputokat.}
\label{fig:arch}
\end{figure}

\subsection{Adatbázis-terv: csillagséma}

Az adatbázis logikai modellje csillagsémára épül, amelyet a 3.\ fejezetben bemutatott általános CDR adatmodellből vezettem le, a YJMob100K sajátosságaihoz igazítva.

A \textbf{ténytáblák} a nyers mobilitási rekordokat tartalmazzák:
\begin{itemize}
    \item \textbf{mobility1} -- a normál időszak 111{,}5~millió rekordja (100~000 felhasználó, 75 nap).
    \item \textbf{mobility2} -- a vészhelyzeti időszak 29{,}4~millió rekordja (25\,000 felhasználó, 75 nap, ebből az utolsó 15 nap rendkívüli). A két adatkészlet független kohorsz, így a teljes felhasználószám $100\,000 + 25\,000 = 125\,000$ (a részleteket lásd az 5.3.\ szakaszban).
\end{itemize}

\noindent Minden rekord öt mezőből áll: \texttt{uid} (előfizető), \texttt{d} (nap), \texttt{t} (30~perces időrés), \texttt{x} és \texttt{y} (rácskoordináták).

A \textbf{dimenzió táblák}:
\begin{itemize}
    \item \textbf{dim\_time} -- az időrések metaadatai: a \texttt{(d, t)} párból számított óra, napszak (éjszaka/reggel/délután/este), hét napja, hétköznap/hétvége jelző, valamint a vészhelyzeti időszak jelölése (d $\geq$ 60). Összesen 3\,600 sor ($75 \times 48$).
    \item \textbf{dim\_cell} -- a rácscellákon számított térbeli attribútumok: domináns POI kategória, zónatípus (élelmiszer, bevásárlás, közlekedés stb.), összesített POI szám. Csak azok a cellák szerepelnek, amelyekben tényleges aktivitás volt (20\,146 sor a 40\,000-ből).
    \item \textbf{dim\_user} -- felhasználói összefoglaló: az adatkészlet azonosítója, becsült otthon- és munkahely-cella, otthon--munkahely távolság. Természetes kulcs: \((\texttt{uid}, \texttt{dataset})\) kompozit kulcs, mivel a két adatkészlet független kohorsz (lásd 5.3.\ szakasz). 121\,607 sor (a 125\,000 felhasználóból ennyihez sikerült stabil otthoncellát hozzárendelni). A felhasználói archetípus-klaszterezés (lásd 9.6.\ szakasz) eredménye nem a dimenzió-táblában, hanem külön \texttt{output/user\_archetypes.csv} fájlban tárolódik, hogy az analitikai modellfuttatások függetlenek maradjanak az adatraktártól.
\end{itemize}

\noindent A származtatott mobilitási metrikákat két aggregált tábla tárolja (\texttt{user\_daily\_features1} és \texttt{user\_daily\_features2}), amelyeket a 4.\ fejezetben bemutatott képletek alapján SQL-ben számítok ki. Ezek felhasználónként és naponként tartalmazzák a gyrációs sugarat, a Shannon-entrópiát, a lokációs diverzitást, a becsült otthon- és munkahely-cellát, valamint az otthon--munkahely távolságot méterben.

\subsection{Feldolgozási folyamat (pipeline)}

A feldolgozást négy, parancssorból egymás után futtatható Python modulra bontom; a köztes adatáramlást a központi \texttt{mobility.duckdb} fájl és Parquet exportok biztosítják. Az egyes fázisok bemenete--kimenete a \autoref{tab:pipeline_phases} táblázatban szerepel; a tényleges implementáció részleteit (kódminták, futási idők) a 7.\ fejezet tárgyalja.

\begin{table}[H]
\centering
\caption{A pipeline négy fázisa, modulok és inputok/outputok.}
\label{tab:pipeline_phases}
\begin{tabular}{l l p{4.5cm} p{4.5cm}}
\hline
\textbf{\#} & \textbf{Modul} & \textbf{Bemenet} & \textbf{Kimenet} \\
\hline
1 & \texttt{ingest.py}   & 4 db CSV (YJMob100K) & 4 DuckDB tábla (\texttt{mobility1/2}, \texttt{cell\_poi\_cat}, \texttt{poi\_categories}) \\
2 & \texttt{features.py} & \texttt{mobility1/2} & \texttt{user\_daily\_features1/2} (DuckDB + Parquet) \\
3 & \texttt{schema.py}   & ténytáblák, POI-tábla & \texttt{dim\_time}, \texttt{dim\_cell}, \texttt{dim\_user} \\
4 & \texttt{anomaly.py}  & \texttt{user\_daily\_features1/2} & \texttt{anomaly\_scores*.parquet} \\
\hline
\end{tabular}
\end{table}

\subsection{Georeferálási megoldás}

Az YJMob100K 200$\times$200-as rácsa önmagában nem hordoz földrajzi információt. A rácsot az EPSG:2449 (Japan Plane Rectangular CS~VII) vetületi koordináta-rendszerbe helyeztem el a \texttt{reverse-engineering-YJMob100K-grid} projekt \cite{pinter2024revealing} eredményei alapján; a vetületi koordináták WGS84-re való átváltását a pyproj \texttt{Transformer} osztálya végzi a \texttt{geo.py} modulban. A rácsparaméterek (origó, anizotrop cellaméretek, tengelykonvenció) és a tényleges implementáció részleteit a 8.\ fejezet tárgyalja.

\subsection{Kockázatok és korlátok}

A rendszerterv megvalósítása során a következő korlátokkal számoltam:
\begin{itemize}
    \item \textbf{Térbeli pontosság:} A rács cellamérete (${\sim}455 \times 556$~m) durvábbá teszi a pozíciómeghatározást, mint egy valódi bázisállomás-hálózat városi részén. A metrikák így a rácsszintű felbontásra vonatkoznak \cite{gonzalez2008understanding}.
    \item \textbf{Georeferálás bizonytalansága:} A rácsparamétereket inverz mérnöki módszerrel határoztam meg; az origó pontos helye néhány száz méteres bizonytalansággal terhelt.
    \item \textbf{Anomáliadetekció validálása:} Felügyelet nélküli módszert alkalmazok (Isolation Forest), amelynek kiértékelése címkézett ground truth adat hiányában korlátozott. Az eredményeket a vészhelyzeti időszak ismert kontextusa alapján értékelem.
\end{itemize}
